\subsection{An exact comparison with the consecutive repeated product}
\label{sec:kt-cubic-shuffle}

Retain all the data and local conditions of
KT-ORDINARY-CUBIC-MASSEY-CONDITION.  All cochains use
\(df(g,h)=f(g)+g f(h)-f(gh)\), and the cup convention is
\((x\smile y)(g,h)=x(g)\,g y(h)\).
Pointwise multiplication of one-cochains is denoted by
juxtaposition.  In particular it is not the cup product.

\begin{proposition}[KT-ORDINARY-CUBIC-SHUFFLE]
Define
\[
 h=-\frac12u^2,\qquad a=-w-\frac12uv,\qquad
 b=w-\frac12uv .                                      \tag{CS.1}
\]
Then \(dh=u\smile u\), \(da=u\smile v\), and
\(db=v\smile u\).  The cocycle
\[
 \Phi=u\smile a+h\smile v\in Z^2(G_{\mathbf Q,S},k(\psi))
                                                               \tag{CS.2}
\]
is a specified consecutive defining value for
\(\langle u,u,v\rangle\).  With
\(\Theta=u\smile b+a\smile u\) as in (CM.20), the exact
cochain identities are
\[
 \Theta+2\Phi=-d(ua),\qquad
 \Omega_+=3\Phi-dR,\qquad
 R=\frac32uw+\frac12u^2v.                              \tag{CS.3}
\]

If \((c,(z_v)_v)\), with \(c=-\Omega_+\), represents
the ordinary obstruction in (CM.14), put
\[
 z^{\mathrm{rep}}_v=\frac{R|_{G_{\mathbf Q_v}}-z_v}{3}.
                                                               \tag{CS.4}
\]
Then \(d z^{\mathrm{rep}}_v=\Phi|_{G_{\mathbf Q_v}}\),
and the relative class
\[
 \mathfrak m_{\mathrm{rep}}(u,u,v)
       =[(\Phi,(z^{\mathrm{rep}}_v)_v)]\in\mathcal O_\psi
                                                               \tag{CS.5}
\]
is well defined under the same permitted global and local
equivalences as the ordinary cubic obstruction.  Exactly
\[
             \mathfrak o_3(u,v)=-3\mathfrak m_{\mathrm{rep}}(u,u,v).
                                                               \tag{CS.6}
\]
Thus changing the order of the repeated entries introduces
no additional gap for this specified relative obstruction.
\end{proposition}

\begin{proof}
The cocycle identity for \(u\) gives
\(d(u^2)=-2u\smile u\).  The other two defining equations
are (CM.19).  The differential and cup rule show that the
two terms in \(d\Phi\) cancel, so \(\Phi\) is a cocycle.

For a direct sign check use the six variables of (CM.9).
They give
\[
\begin{split}
 \Phi(g,h)&=-Ac-\frac12Aab-\frac12A^2b,\\
 \Theta(g,h)&=Ac-Ca-\frac12Aab-\frac12ABa,\\
 d(uw)(g,h)&=-Ac-Ca-\frac12(A+a)(Ab-Ba),\\
 d(u^2v)(g,h)&=-A^2b-2ABa-2Aab-Ba^2.
\end{split}                                                   \tag{CS.7}
\]
Substitution proves the first identity in (CS.3), since
\(-ua=uw+u^2v/2\).  Substituting that identity in (CM.21)
gives the second identity, with precisely the displayed \(R\).
All denominators are units because \(p>3\).

For the relative assertion, (CS.3) gives
\(c=-3\Phi+dR\), and (CS.4) is therefore a local primitive
of \(\Phi\).  At the level of pairs one has the literal equality
\[
 (c,(z_v)_v)=-3(\Phi,(z^{\mathrm{rep}}_v)_v)
                              +(dR,(R|_v)_v).                 \tag{CS.8}
\]
The last pair is zero in (CM.14).  Since the allowed local
spaces \(L_v\) are vector subspaces and \(3\) is invertible,
any change by an allowed local cocycle transports to an
allowed local cocycle.  A change by a global primitive
transports to the corresponding global primitive.
The already proved independence of \(\mathfrak o_3\),
together with (CS.8), proves the asserted independence
of (CS.5).  This transports the actual ordinary local
choices; it does not replace them by arbitrary local
primitives for an unrestricted Massey product.
\end{proof}

\paragraph{The exact relation to Sharifi's conventions.}
Here is a further algebraic comparison, separate from the
proper norm-system requirement in Sharifi's theorem.
Restrict to \(K=\mathbf Q(\sqrt D,\mu_p)\), so the
coefficient characters become trivial, and first take
\(k=\mathbf F_p\).  A fixed primitive \(p\)-th root
identifies the coefficients with \(\mu_p\).
In the upper-unitriangular four-by-four group modulo its
central \((1,4)\) entry, the functions
\[
 \kappa_{12}=\kappa_{23}=u,\quad\kappa_{34}=v,\quad
 \kappa_{13}=-h=\frac12u^2,\quad\kappa_{24}=-a
                                                               \tag{CS.9}
\]
give a homomorphism.  Indeed Sharifi's defining equation
(his condition (c), printed p.~18) is
\(d\kappa_{ij}=-\sum_l\kappa_{il}\smile\kappa_{lj}\),
which is exactly (CS.1) with these signs.
Conjugation by \(I-E_{12}/2\) changes only the noncentral
entry \(\kappa_{13}\), replacing it by
\(\binom{u}{2}=(u^2-u)/2\).  Its values therefore lie
in the subgroup \(G_1^{(2)}/Z_1^{(2)}\) generated in
Sharifi's notation by the images of
\(X=I+E_{12}+E_{23}\) and \(Y=I+E_{34}\).
There is no remaining discrepancy in the repeated-entry
block or its binomial normalization.

With the convention (4.1) of that paper, the cocycle
of this last defining system is
\[
 \kappa_{12}\smile\kappa_{24}
       +\kappa_{13}\smile\kappa_{34}
       =-\Phi-\frac12d a .                              \tag{CS.10}
\]
Thus its raw class is \(-[\Phi]\), and the raw ordinary
cubic obstruction is three times this class.  The same
global-coboundary transport as (CS.8) adjusts the local
primitives.  For larger finite coefficient fields the
cochain identities remain valid over that field; a
literal application of a theorem with \(\mu_p\)
coefficients must still identify the requisite scalar
components rather than treating a general \(k\)-valued
character as one cyclic degree-\(p\) extension.

\begin{corollary}[KT-CUBIC-BOCKSTEIN-NORM-IDENTIFICATION]
Suppose \(k=\mathbf F_p\) and \([u]\ne0\).
Over \(K=\mathbf Q(\sqrt D,\mu_p)\), let \(S_K\) be
the primes above \(S\), set \(G=G_{K,S_K}\), and write
\(\chi=u|_G\).  This is a surjection onto \(H=\mathbf F_p\).
Let \(\mathcal R=\mathbf F_p[H]\), choose \(\tau\in H\)
corresponding to \(1\in\mathbf F_p\), and put \(x=[\tau]-1\),
\(I=(x)\).  The latter notation means the element of the
coimage corresponding to \(1\); it does not require a
finite-order lift in \(G\).
Then
\[
 f=v-a x\in Z^1(G,\mathcal R/I^2),\qquad
 \Psi^{(2)}([f])=-[\Phi|_G]x^2,                       \tag{CS.11}
\]
where \(G\) acts on \(\mathcal R\) by multiplication
by its image in \(H\), and \(\Psi^{(2)}\) is the
connecting map for the augmentation sequence through
\(I^3\).  The conjugated system (CS.9) is a binomial proper
system in the sense of Lam--Liu--Sharifi--Wake--Wang.
Its realization by a Sharifi norm system additionally
uses its zero absolute obstruction and
KT-CUBIC-MATCHED-NORM-SYSTEM below.

In particular, choose a primitive root \(\zeta\in\mu_p\)
and Kummer representatives \(A_0,B_0\in K^\times\)
for \(\chi,v|_G\) using \(\zeta\).  Put
\(L=K(A_0^{1/p})\), of degree \(p\), and let
\(\Omega_S\) be the maximal extension of \(K\)
unramified outside \(S_K\).  The matched-system
argument supplies a norm-system representative
\(y\in L^\times\) with
\[
 N_{L/K}y=B_0,\qquad D^{(1)}y\in\Omega_S^{\times p}.
                                                               \tag{CS.12}
\]
Sharifi's Theorem~5.4 applies to this system.  If
\(\mathfrak Y\) is the ideal appearing there and
\(P^{(1)}_{1,L/K}\) is its lower-product subgroup,
then the raw image \(\overline{\mathfrak o}_3\) of
the relative ordinary obstruction satisfies
\[
 \iota_\zeta\!\left(\operatorname{res}_{K/\mathbf Q}
                     \overline{\mathfrak o}_3\right)
     \equiv -3[N_{L/K}\mathfrak Y]
             \pmod{P^{(1)}_{1,L/K}}.                  \tag{CS.13}
\]
Here \(\iota_\zeta:\mathbf F_p\to\mu_p\) sends
\(1\) to \(\zeta\); equivalently one untwists
Sharifi's triple coefficient by \(\zeta^{\otimes2}\).
A nonzero right-hand coset in (CS.13) is sufficient
for a nonzero ordinary relative cubic obstruction.
No such nonzero coset is asserted here.
\end{corollary}

\begin{proof}
Restriction is injective in degree one because
\([K:\mathbf Q]\) is prime to \(p\).  The nonzero
\(\mathbf F_p\)-valued character \(\chi\) is therefore
surjective.  The action of \(g\in G\) on the group
ring is multiplication by \((1+x)^{\chi(g)}\).
The equations \(dv=0\), \(da=u\smile v\) show
directly that \(f\) is a cocycle modulo \(x^2\).
Use the same expression as a one-cochain modulo \(x^3\).
Its differential has coefficient of \(x^2\) equal to
\[
 u(g)(-a(h))+\binom{u(g)}2v(h),
                                                               \tag{CS.14}
\]
which is precisely (CS.10).  This proves (CS.11),
including its sign, without any comparison of
cohomology dimensions.

Definition~3.3.2 and the complete Lemma~3.3.3 of
Lam--Liu--Sharifi--Wake--Wang identify the binomial
block in (CS.9) as a proper partial-system extension.
In their Section~4.3 take \(n=2\), \(A=R=\mathbf F_p\):
the size condition is \(2p<p^2\), and \(p>3\)
ensures it.  Their Theorem~4.3.1 gives exactly the
same Bockstein class and fixes the binomial matrix
conventions.  That terminology alone does not
realize every specified last column by a norm
representative.  Here the additional hypothesis
needed by KT-CUBIC-MATCHED-NORM-SYSTEM holds:
the ordinary obstruction is locally zero everywhere
by (CM.14), including unramified places outside
\(S\).  After restriction to \(K\), its inflation
to \(H^2(G_K,\mu_p)=\operatorname{Br}(K)[p]\)
is zero by the Brauer local-global injection.
By (CS.10) its value is three times the Bockstein,
so the latter also has zero absolute inflation.
The matched-system lemma lifts the specified
truncated cocycle through the absolute regular
module, then uses Shapiro and the exact Kummer-dual
pairing to give (CS.12).  This argument supplies
the norm representative for the actual defining
class.  The lemma itself is independent of the
norm formula being applied here.

Apply Sharifi's Theorem~5.4 with his \(m=n=1\)
and repeated length \(k=2\).  Since \(2<p-1\),
its exceptional boundary term is absent.  Its
formula for the proper triple value is the
negative norm-ideal class.  Equations (CS.6)
and (CS.10) make the raw ordinary value three
times that triple value, proving (CS.13).
If its coset is nonzero then the restricted raw
class is nonzero, so the relative class cannot
vanish.  This establishes the final sufficient
condition without asserting that it is met.
\end{proof}

\paragraph{The remaining arithmetic requirements.}
The \(\mathbf F_p\)-valued ordering and norm-system
issues are discharged by (CS.3), (CS.9), the zero
absolute inflation, and the explicit matched-system
lemma.  The raw formula (CS.13) by itself does not
determine the local correction in (CM.15).
KT-CUBIC-EXCISION and KT-CUBIC-P-ONLY-NORM below
retain that correction in a class with ramification
only above \(p\), and give an exact ideal-class
formula when the norm representative is matched.
No norm-ideal class has been evaluated here;
neither nonvanishing nor injectivity in the
original strict class has been proved.
General finite \(k\)-valued input still
requires a scalar or cyclic reduction to use this
particular norm formula.  These are equivalences
and an available sufficient evaluation criterion,
not a new arithmetic exclusion or a source-to-jet
theorem.

The primary passages inspected for this comparison are
Sharifi, \emph{Massey products and ideal class groups},
author version dated 26 July 2006: Section~4, printed
pp.~17--22, especially condition (c), (4.1), the two
generators, Proposition~4.2 and the proper-system
definition; and Theorem~5.4 with its complete proof,
printed pp.~26--27.  The published article is
\emph{J.~reine angew.~Math.}~603 (2007), 1--33.
For the further imported properness identification
we read Lam--Liu--Sharifi--Wake--Wang,
\emph{Generalized Bockstein maps and Massey products},
\emph{Forum Math.~Sigma}~11 (2023), e5:
Definition~3.3.2, the complete proofs of
Lemma~3.3.3 and Theorem~3.3.4, and the complete
Section~4.3 (author-PDF pp.~24--27 and 30--32).
The identification paragraph and Theorem~4.3.1
were also checked in the published Cambridge
version.  Its Section~6 absolute-Galois vanishing
statements do not imply vanishing for this fixed
ramification set with the prescribed local
primitives.
